% ===================================================================
%  Chemie: Energiediagramm einer exothermen Reaktion, mit und ohne
%  Katalysator. Die Kurven sind Funktionen – Energien oben ändern.
%  Kompilieren: pdflatex vorlage_energiediagramm.tex
% ===================================================================
\documentclass[border=2mm]{standalone}
\usepackage[T1]{fontenc}
\usepackage{lmodern}
\usepackage[ngerman]{babel}
\usepackage{amsmath}
\usepackage{amssymb}
\usepackage{tikz}
\usetikzlibrary{arrows.meta}

\definecolor{akzentblau}{HTML}{2563EB}
\definecolor{akzentrot}{HTML}{DC2626}
\definecolor{akzentgruen}{HTML}{16A34A}

% ---- Werte (in cm auf dem Papier) ---------------------------------
\def\Eedukt{3}       % Energie der Edukte
\def\Eprodukt{1}     % Energie der Produkte (kleiner -> exotherm)
\def\Ea{2.6}         % Aktivierungsenergie ohne Katalysator
\def\Eak{1.2}        % Aktivierungsenergie mit Katalysator
% --------------------------------------------------------------------
% Verlauf: weicher Übergang (Sigmoid) von Edukt- zu Produktniveau
% plus ein "Berg" (Gaußglocke) in der Mitte
\begin{document}
\begin{tikzpicture}[>={Stealth[length=5pt]}, font=\small,
  declare function={
    stufe(\x) = \Eedukt + (\Eprodukt-\Eedukt)/(1+exp(-2.2*(\x-4)));
    berg(\x,\h) = stufe(\x) + (\h + (\Eedukt-stufe(4)))*exp(-(\x-4)^2/0.9);}]
  \draw[->] (0,0) -- (8.4,0) node[below left] {Reaktionsverlauf};
  \draw[->] (0,0) -- (0,6.4) node[above] {Energie};
  \draw[akzentrot, very thick, smooth, samples=120, domain=0.3:7.8] plot (\x, {berg(\x,\Ea)});
  \draw[akzentgruen!70!black, very thick, dashed, smooth, samples=120, domain=0.3:7.8]
       plot (\x, {berg(\x,\Eak)});
  % Niveaus und Pfeile
  \draw[dotted] (0.3,\Eedukt) -- (5.2,\Eedukt);
  \draw[dotted] (4.3,\Eprodukt) -- (7.8,\Eprodukt);
  \draw[dotted, akzentrot] (2.6,\Eedukt+\Ea) -- (4,\Eedukt+\Ea);
  \draw[dotted, akzentgruen!70!black] (3.3,\Eedukt+\Eak) -- (4,\Eedukt+\Eak);
  \draw[<->, akzentrot] (2.6,\Eedukt) -- node[left] {$E_A$} (2.6,\Eedukt+\Ea);
  \draw[<->, akzentgruen!70!black] (3.3,\Eedukt) -- node[right, xshift=-1pt] {$E_A'$} (3.3,\Eedukt+\Eak);
  \draw[<->, akzentblau] (6.8,\Eprodukt) -- node[right] {$\Delta_r H < 0$} (6.8,\Eedukt);
  \draw[dotted] (6.8,\Eedukt) -- (5.2,\Eedukt);
  \node[above] at (0.9,\Eedukt) {Edukte};
  \node[below] at (7.2,\Eprodukt) {Produkte};
  \draw[akzentrot, very thick] (5.6,5.9) -- (6.1,5.9) node[right, text=black] {ohne Katalysator};
  \draw[akzentgruen!70!black, very thick, dashed] (5.6,5.4) -- (6.1,5.4) node[right, text=black] {mit Katalysator};
\end{tikzpicture}
\end{document}
